\icode{E-MI-SI-L3-261106-003}
夏日天气十分炎热，为了提高空调的制冷效果，门窗都需紧闭。某学校教学楼有些门因为老化故障无法紧闭，细心的小明给门下缝隙处塞紧了一个小木楔，侧面如图所示。已知木楔质量为\(m\)，其上表面可视作光滑，下表面与水平地面间的动摩擦因数为\(\mu\)，木楔上表面与水平地面间夹角为\(\theta\)，重力加速度为\(g\)，最大静摩擦力等于滑动摩擦力，不计木楔的形变。若有人轻推门时，门给木楔一个大小为\(F\)的力推动木楔在地板上缓慢匀速移动，下列说法正确的是\paren[]

\begin{tikzpicture}
	\draw (0,0)--(15:1.2)|-(0,0);
	\draw[thin,fill=orange!60] (15:.8)coordinate(A) rectangle($(A)+(-.2,2.5)$);
	\foreach \x in{-.4,-.2,...,1.4}
		\draw[gray] (\x,-.2)--++(.2,.2);
	\draw[thick] (-.4,0)--(1.6,0);
\end{tikzpicture}
\begin{choices}
	\item 力\(F\)沿水平方向的分力大小为\(F \cos \theta\)
	\item 木楔受到的滑动摩擦力大小为\(\mu mg\)
	\item 力\(F\)的大小为\(\frac{\mu mg}{\sin \theta - \mu \cos \theta}\)
	\item 若\(\tan \theta_0 = \mu\)，木楔倾角\(\theta \geq \theta_0\)时，不管用多大力推门，木楔都能将门卡住而不再运动
\end{choices}
\begin{answer}
    
\end{answer}